% ---------------------------------------------------------------------------
% Capítulo 3 --- Metodología
% ---------------------------------------------------------------------------
% Mínimo dos páginas por capítulo (norma 2.6). Declara las decisiones de
% diseño con su fecha, no solo describe el sistema resultante: son las que
% hacen reproducible el experimento del capítulo 4.

\chapter{Metodología}
\label{cap:metodologia}

El trabajo se organizó en seis fases sucesivas, correspondientes una a una con los objetivos
específicos enunciados en la Introducción, más una séptima fase de consolidación documental. Cada
fase produjo un entregable concreto sobre el que se apoya la siguiente, y las decisiones de diseño
que no podían anticiparse al inicio ---porque dependían de una medición, no de un criterio--- se
tomaron y registraron en el momento en que la medición estuvo disponible, con su fecha. Esta
sección documenta esas decisiones en el orden en que se tomaron.

% ===========================================================================
\section{Enfoque general: modelar en lugar de esperar}
% ===========================================================================

La primera decisión metodológica, y la que condiciona a todas las demás, se tomó antes de escribir
ninguna línea de código: al no disponer de acceso a un cliente real ni a su sistema propietario de
análisis de alertas, el trabajo optó por \emph{modelar} un cliente genérico en lugar de
\emph{analizar} uno concreto. Analizar un cliente real habría producido un prototipo ajustado a ese
cliente; modelar uno genérico obligó a separar explícitamente, desde el principio, lo que cualquier
empresa comparte ---una infraestructura de red organizada por capas, con una superficie de gestión
remota presente en todas ellas--- de lo que varía de una a otra ---su catálogo de acciones permitidas,
su política de continuidad de servicio, la criticidad de sus activos---. Esa separación es la que
permite que el prototipo resultante sea configurable por cliente sin cambios de código, y es la
misma lógica que, más adelante, resolvió la ausencia de una fuente real de alertas: en vez de
esperar a que un sistema propietario estuviera disponible, se sustituyó por una plataforma de
detección basada en reglas ---Wazuh--- desplegada dentro de un laboratorio propio, con la ventaja
adicional de que su nivel de regla podía servir directamente como método tradicional de referencia.

% ===========================================================================
\section{Fase 1: modelado del cliente y del caso de uso}
% ===========================================================================

Se construyó un modelo de cliente genérico describiendo su infraestructura de red por capas ---punto
de entrega, equipo de borde, conmutación, zona desmilitarizada, servidores internos, puestos de
trabajo, dispositivos IoT y plano de gestión---, identificando en cada una qué expone hacia la
administración remota y qué ocurre si se interrumpe. De ese modelo se derivaron cinco puntos de
variabilidad configurables entre clientes ---el esquema de la alerta, el repertorio de acciones y su
canal, la política de continuidad, el inventario y criticidad de activos, y la severidad de
referencia--- y un conjunto de restricciones operativas que no son negociables en ningún cliente:
que ningún servicio se detenga sin autorización, que no exista una ventana de mantenimiento sobre
la que apoyarse, que toda acción sea reversible porque no hay manos físicas junto al equipo, y que
ninguna acción pueda cortar el propio canal de gestión por el que se administra y se responde.

Sobre ese modelo se acotó el caso de uso a una sola familia de alertas ---el acceso no autorizado a
servicios de gestión expuestos---, deliberadamente definida por el tipo de patrón y no por el equipo
concreto que lo sufre: la misma fuerza bruta contra un servicio de administración se reconoce igual
en un conmutador que en un servidor de ficheros, lo que permite que el sistema generalice al parque
de equipos de cualquier cliente en lugar de aprender las particularidades de uno solo.

% ===========================================================================
\section{Fase 2: estado del arte y registro de requisitos}
% ===========================================================================

Se revisaron, por un lado, las plataformas comerciales de detección y respuesta extendida y
gestionada y sus limitaciones documentadas, y por otro, las aplicaciones de modelos de lenguaje a
la seguridad, incluyendo sus riesgos específicos. De ese análisis se derivó un registro único de
veinte requisitos funcionales y catorce no funcionales, cada uno trazado hasta la limitación de
mercado o el criterio de diseño que lo motiva. Cuando, más adelante, el encuadre del proyecto se
desplazó del segmento de pequeña y mediana empresa hacia la red de un cliente modelado en la Fase 1,
ese registro se revisó íntegramente contra el nuevo contexto: ningún requisito original dejó de
aplicar, seis cambiaron de matiz y siete se añadieron ---principalmente los que más adelante
formalizarían la política de continuidad de servicio---, lo que se documentó como una revisión
explícita y no como una reescritura silenciosa.

% ===========================================================================
\section{Fase 3: construcción del entorno de pruebas}
% ===========================================================================

El entorno de pruebas se construyó como una red de cliente emulada con Containerlab, reproduciendo
desde el punto de entrega del proveedor hacia dentro: un equipo de borde, un conmutador de red
interna, un puesto de usuario, un dispositivo de tipo IoT con servicios expuestos y un servidor con
vulnerabilidades documentadas ---Metasploitable2, con diecinueve puertos abiertos y fallos conocidos
y públicos---, todos conectados además a una red de gestión fuera de banda por la que operan, sin
contaminar el tráfico observado, el motor de triaje y un nodo auditor con Nmap.

Un mismo laboratorio debía resolver dos necesidades que, se verificó explícitamente, no son
intercambiables: un \emph{ground truth} de vulnerabilidades ---qué falla exactamente en cada
nodo--- y un \emph{ground truth} de alertas ---si una alerta concreta es un verdadero o un falso
positivo---. El primero sale directamente de la composición documentada de la topología; el segundo
exigió un puente explícito: cada alerta que Wazuh generaba sobre la actividad del laboratorio se
contrastaba contra el inventario de vulnerabilidades del nodo al que apuntaba, y de esa
comparación ---con revisión manual, no completamente automatizable--- salía la etiqueta de verdadero o
falso positivo, conservando en todo momento el nivel de regla original de Wazuh como el dato que
más adelante serviría de baseline.

El resultado, cerrado el 31 de agosto de 2026, fue un conjunto de 410 alertas reales generadas
sobre el propio laboratorio, etiquetadas y particionadas por campaña: 20 verdaderos positivos, 16
falsos positivos y 374 alertas de ruido de plataforma ---telemetría de evaluación de configuración
que no corresponde a ningún patrón del caso de uso acotado---, todas ellas concentradas, las
soportadas, en un único activo y servicio (acceso por SSH al servidor con vulnerabilidades
documentadas). Esta concentración en una sola familia de ataque, verificada al construir el
conjunto y no descubierta después, es la que más adelante bloquearía el ajuste fino del
clasificador y limitaría el alcance estadístico de la evaluación del capítulo \ref{cap:resultados}:
se documenta aquí como una condición conocida del método, no como un hallazgo posterior.

% ===========================================================================
\section{Fase 4: diseño de la arquitectura}
% ===========================================================================

El diseño se apoyó en tres puntos de aislamiento deliberados, cada uno resuelto mediante un
adaptador que esconde una dependencia externa hoy provisional: un módulo de ingesta que aísla la
fuente de alertas concreta ---hoy Wazuh, mañana el sistema propietario del cliente---; una interfaz de
análisis de dos operaciones, \texttt{clasificar} y \texttt{justificar}, que aísla el modelo de
lenguaje concreto detrás de cada una; y un conector que aísla el protocolo de comunicación hacia el
equipo objetivo ---hoy SSH sobre un catálogo cerrado de acciones---. Ninguno de los tres exige
reescribir la lógica de decisión del motor de triaje cuando la pieza que ocultan cambie.

Sobre esa interfaz de análisis se definieron dos perfiles de despliegue según el hardware
disponible: un perfil híbrido para el equipo de desarrollo entonces disponible ---un procesador de
propósito general de ocho núcleos, sin GPU dedicada y con dieciséis gigabytes de memoria---, que
combina un codificador ligero para clasificar con un modelo generativo pequeño para justificar en
línea; y un perfil de un único componente generativo especializado, viable únicamente con una GPU
dedicada o memoria sustancialmente mayor. La decisión inicial del 24 de agosto de 2026 estimaba,
sobre mediciones publicadas en hardware comparable, que un modelo generativo de tres mil ochocientos
millones de parámetros cuantizado rendiría entre diez y doce \emph{tokens} por segundo en el equipo
disponible. La medición directa realizada el 2 de septiembre de 2026 dio un resultado muy distinto
---aproximadamente 3,7 \emph{tokens} por segundo---, lo que se documentó de inmediato como una
corrección y no se ocultó: con ese rendimiento, un modelo de ese tamaño quedaba fuera del
presupuesto de latencia tolerable para una validación humana interactiva, y el justificador se
implementó finalmente sobre un modelo generativo de mil millones de parámetros. La interfaz de
análisis, precisamente por estar aislada del modelo concreto, no tuvo que rediseñarse para absorber
ese cambio.

También en la Fase 4 se resolvió que la decisión de qué acción proponer nunca recae en el modelo de
lenguaje, por las razones expuestas en la sección \ref{sec:politica-perfil} del marco teórico: una
tabla de política determinista, indexada por la clase de la alerta, propone la acción de menor
impacto que resuelve la amenaza, y un perfil de cliente configurable la filtra ---permitiéndola,
degradándola a una de menor impacto o reteniéndola para validación humana--- según lo que ese cliente
concreto permite. Se definió además el plan de métricas de la evaluación ---precisión, exhaustividad,
medida $F_1$, tasa de falsos positivos, acierto de prioridad y dos métricas de continuidad
específicas de este caso de uso: acciones disruptivas indebidas y retención correcta---, y se fijó
como método tradicional de referencia el nivel de regla de Wazuh, por ser literalmente el enfoque
basado en reglas y firmas contra el que el proyecto se planteó comparar.

% ===========================================================================
\section{Fase 5: implementación}
% ===========================================================================

El prototipo se implementó en Python 3, empleando únicamente su biblioteca estándar más
\texttt{pyyaml} para la configuración, sin gestor de paquetes ni entorno virtual adicional para el
motor de decisión, y con \texttt{unittest} como marco de pruebas. La implementación se dividió en
tres subproyectos:

\begin{itemize}
  \item \textbf{Núcleo de decisión.} La cadena ingesta~$\rightarrow$~análisis~$\rightarrow$~política~
        $\rightarrow$~perfil~$\rightarrow$~traza, ejecutable en lote sobre el conjunto de alertas
        etiquetado, con un clasificador de referencia determinista ---una regla fija sobre la postura
        de exposición del activo, sin aprendizaje--- detrás de la interfaz \texttt{clasificar}. Este
        clasificador se mantiene como una pieza provisional deliberada: el ajuste fino de un
        codificador quedó bloqueado precisamente por la concentración en una sola familia de ataque
        documentada en la Fase 3, sin variedad suficiente de clases para entrenar y validar un
        modelo que generalizara.
  \item \textbf{Lazo en vivo.} Un conector que traduce cada acción abstracta del catálogo cerrado a
        un comando concreto sobre el nodo objetivo por SSH, con verificación de idempotencia previa
        a la ejecución, validación de los parámetros contra caracteres de control de shell, y una
        segunda verificación posterior para confirmar el efecto; una validación humana por terminal
        que muestra la alerta, la clase, la prioridad, la justificación y la acción propuesta antes
        de ejecutarla; y el registro de traza auditable de cada decisión. El ejecutor SSH real se
        inyecta como una función sustituible, de forma que la lógica de idempotencia y de
        verificación se pudo probar con un ejecutor simulado sin necesidad del laboratorio activo.
  \item \textbf{Justificador con modelo de lenguaje real.} Un modelo generativo de mil millones de
        parámetros, cuantizado a cuatro bits, invocado mediante un proceso independiente a un motor
        de inferencia en un entorno Conda separado del resto del prototipo. Cada justificación se
        somete a una verificación de anclaje automática ---ninguna dirección IP mencionada puede ser
        distinta de la de la alerta, y el texto debe referenciar al menos un campo concreto de
        esta--- y, si no la supera o el modelo falla, se degrada a una plantilla determinista que por
        construcción sí ancla, de forma que ninguna alerta queda nunca sin justificación.
\end{itemize}

Sobre esa misma interfaz se añadió un cuarto componente, no previsto en el diseño original de la
Fase 4 sino incorporado tras detectar en la Fase 6 que el justificador de mil millones de
parámetros, aunque anclaba correctamente sus afirmaciones, las describía con un significado
equivocado: un módulo de generación aumentada por recuperación que indexa, mediante un modelo de
\emph{embeddings} del mismo entorno de inferencia, un corpus curado de fichas sobre las técnicas
MITRE ATT\&CK, las reglas de Wazuh y las vulnerabilidades del laboratorio pertinentes al caso de
uso, y recupera los tres pasajes más afines a cada alerta ---por similitud coseno, según la ecuación
de la sección \ref{sec:problema-triaje-referencia} del marco teórico--- antes de invocar al modelo
generativo. Dos decisiones de diseño acotan deliberadamente su alcance: la consulta de recuperación
se construye únicamente a partir de campos estructurados de la alerta ---el identificador de la
regla, la técnica y el servicio---, nunca del registro completo del evento, por la misma razón de
seguridad que ya impedía tratar ese campo como instrucción; y el corpus se mantiene reducido y
curado a mano, en vez de indexar documentación completa, porque el objetivo era verificar si
corregir el conocimiento del modelo mejoraba la corrección de su justificación, no maximizar la
cobertura de la recuperación.

Todo el prototipo se validó con una batería de pruebas unitarias que, en el momento de escribir
este informe, asciende a más de un centenar de casos, ejecutados con generadores y ejecutores
simulados que permiten probar la lógica de decisión, de idempotencia y de degradación sin necesidad
de un modelo de lenguaje cargado ni de un laboratorio en ejecución.

% ===========================================================================
\section{Fase 6: evaluación}
% ===========================================================================

La evaluación se diseñó como una campaña reproducible por línea de comandos que procesa una
partición del conjunto de alertas etiquetado, calcula la clasificación del prototipo y, sobre la
misma partición, calcula el mejor punto de un barrido del nivel de regla de Wazuh como umbral de
clasificación ---el método tradicional de referencia---, de forma que ambos se comparan siempre bajo
las mismas condiciones y contra el mismo \emph{ground truth}. Para la calidad de la justificación se
definieron dos evaluaciones complementarias: una automática, que mide la tasa de anclaje ---el
porcentaje de justificaciones cuyas afirmaciones referencian campos que existen realmente en la
alerta---, y una manual, sobre una muestra, en la que un revisor humano juzga si el contenido de la
explicación es correcto y no solo si cita los campos adecuados. Esta distinción resultó decisiva: es
precisamente la que permitió detectar que un anclaje perfecto no garantiza una justificación
semánticamente correcta, y la que motivó incorporar la recuperación aumentada descrita en la
sección anterior.

Para aislar el efecto de esa incorporación, la campaña se ejecutó dos veces sobre el mismo
subconjunto de dieciocho alertas soportadas de la partición de evaluación, con el modelo generativo
y la temperatura de muestreo fijadas en ambos casos: una vez con el justificador simple, y otra con
la recuperación aumentada activada, comparando en cada caso el mismo par de alertas para poder
atribuir cualquier diferencia observada a la sola presencia del módulo de recuperación.
